% GENERATED FILE. Edit canonical lessons, then run npm run generate:books.
% canonical-source-hash: fnv1a64:ab33dbeba6685489
% canonical-lessons: MW-C228-lukho, MW-C228-tapno, MW-C228-pasino, MW-C228-silo, MW-C228-sukno, MW-R228-review-223-228

\chapter{Summer Sun}
\label{ch:mw-summer-sun}
\hlchaptermodality{228}

\begin{chapteropening}
\textbf{By the end of this chapter:} \emph{I can describe dry air, the heat and a cool evening.}

Complete the last lesson of chapter 228: I can describe dry air, the heat and a cool evening.
\end{chapteropening}

\section[\texorpdfstring{\mw{लूखो} (lūkho)}{लूखो (lūkho)}]{\mw{लूखो} (lūkho) — parched; dry and rough}

\label{lesson:MW-C228-lukho}



\begin{quote}
Before the new one: say the Marwadi for a Persian wheel, then the Marwadi for a drought.
\end{quote}

\subsection*{You'll want to know: \mw{लूखो}}
\textbf{\mw{लूखो}} — \emph{lūkho} — \textquotedblleft{}parched; dry and rough\textquotedblright{}.

\begin{quote}
\mw{आज} \mw{हवा} \mw{लूखी} \mw{है।}
\end{quote}

\begin{quote}
\emph{āj havā lūkhī hai.}
\end{quote}

\begin{quote}
The air is parched today.
\end{quote}

\begin{cousinweb}
\textbf{\mw{लूखो}} (\emph{lūkho}) is inherited from Sanskrit \emph{rūkṣa}, \textquotedblleft{}dry, rough\textquotedblright{}; CDIAL lists Hindi \emph{lūkhā}, \textquotedblleft{}dry, not greasy\textquotedblright{}, and a Gujarati cousin. CDIAL gives no Marwari form in this entry (CDIAL 10799).
\end{cousinweb}

\begin{grammarlens}[title={What you've built}]
One more word for this chapter.
\end{grammarlens}

\subsection*{Your turn}
\begin{itemize}
\raggedright
  \item \emph{Say it:} \emph{lūkho}
  \item \emph{Say it:} \emph{lūkho}, once more
  \item \emph{Say it:} \emph{āj havā lūkhī hai.} — The air is parched today.
  \item \emph{Recall:} say the Marwadi for a Persian wheel, then the Marwadi for a drought, then say \emph{lūkho} again
\end{itemize}

\begin{tcolorbox}[breakable,colback=teal!4,colframe=teal!35!black,title={Before you move on}]
What does \mw{लूखो} mean? (\textquotedblleft{}Parched; dry and rough\textquotedblright{}.) Say it once more.
\end{tcolorbox}
\section[\texorpdfstring{\mw{तपणो} (tapṇo)}{तपणो (tapṇo)}]{\mw{तपणो} (tapṇo) — to blaze; to be hot (of the sun)}

\label{lesson:MW-C228-tapno}



\begin{quote}
Before the new one: say the Marwadi for a drought, then the Marwadi for parched.
\end{quote}

\subsection*{You'll want to know: \mw{तपणो}}
\textbf{\mw{तपणो}} — \emph{tapṇo} — \textquotedblleft{}to blaze; to be hot (of the sun)\textquotedblright{}.

\begin{quote}
\mw{जेठ} \mw{मांय} \mw{सूरज} \mw{घणो} \mw{तपै} \mw{है।}
\end{quote}

\begin{quote}
\emph{jeṭh mā̃y sūraj ghaṇo tapai hai.}
\end{quote}

\begin{quote}
In early summer the sun blazes.
\end{quote}

Past (it …): \mw{वो} \mw{सूरज} \mw{तप्यो}, \emph{vo sūraj tapyo}.

\begin{cousinweb}
\textbf{\mw{तपणो}} (\emph{tapṇo}) is inherited from Sanskrit \emph{tapati}, \textquotedblleft{}is hot, heats\textquotedblright{}; CDIAL lists Hindi and Gujarati verbs for being heated under it, and Marwadi \emph{tapṇo} keeps the old \emph{p}. CDIAL gives no Marwari form in this entry (CDIAL 5671).
\end{cousinweb}

\begin{grammarlens}[title={What you've built}]
One more word for this chapter.
\end{grammarlens}

\subsection*{Your turn}
\begin{itemize}
\raggedright
  \item \emph{Say it:} \emph{tapṇo}
  \item \emph{Say it:} \emph{tapṇo}, once more
  \item \emph{Say it:} \emph{jeṭh mā̃y sūraj ghaṇo tapai hai.} — In early summer the sun blazes.
  \item \emph{Recall:} say the Marwadi for a drought, then the Marwadi for parched, then say \emph{tapṇo} again
\end{itemize}

\begin{tcolorbox}[breakable,colback=teal!4,colframe=teal!35!black,title={Before you move on}]
What does \mw{तपणो} mean? (\textquotedblleft{}To blaze; to be hot (of the sun)\textquotedblright{}.) Say it once more.
\end{tcolorbox}
\section[\texorpdfstring{\mw{पसीनो} (pasīno)}{पसीनो (pasīno)}]{\mw{पसीनो} (pasīno) — sweat}

\label{lesson:MW-C228-pasino}



\begin{quote}
Before the new one: say the Marwadi for parched, then the Marwadi for to blaze.
\end{quote}

\subsection*{You'll want to know: \mw{पसीनो}}
\textbf{\mw{पसीनो}} — \emph{pasīno} — \textquotedblleft{}sweat\textquotedblright{}.

\begin{quote}
\mw{जेठ} \mw{मांय} \mw{पसीनो} \mw{आवै} \mw{है।}
\end{quote}

\begin{quote}
\emph{jeṭh mā̃y pasīno āvai hai.}
\end{quote}

\begin{quote}
You sweat in early summer.
\end{quote}

A masculine noun.

\begin{cousinweb}
\textbf{\mw{पसीनो}} (\emph{pasīno}) is inherited from Sanskrit \emph{prasvinna}, \textquotedblleft{}sweated\textquotedblright{}; CDIAL lists Hindi \emph{pasīnā} and a Gujarati word for sweat. CDIAL gives no Marwari form in this entry (CDIAL 8897).
\end{cousinweb}

\begin{grammarlens}[title={What you've built}]
One more word for this chapter.
\end{grammarlens}

\subsection*{Your turn}
\begin{itemize}
\raggedright
  \item \emph{Say it:} \emph{pasīno}
  \item \emph{Say it:} \emph{pasīno}, once more
  \item \emph{Say it:} \emph{jeṭh mā̃y pasīno āvai hai.} — You sweat in early summer.
  \item \emph{Recall:} say the Marwadi for parched, then the Marwadi for to blaze, then say \emph{pasīno} again
\end{itemize}

\begin{tcolorbox}[breakable,colback=teal!4,colframe=teal!35!black,title={Before you move on}]
What does \mw{पसीनो} mean? (\textquotedblleft{}Sweat\textquotedblright{}.) Say it once more.
\end{tcolorbox}
\section[\texorpdfstring{\mw{सीळो} (sīḷo)}{सीळो (sīḷo)}]{\mw{सीळो} (sīḷo) — cool; pleasantly cold}

\label{lesson:MW-C228-silo}



\begin{quote}
Before the new one: say the Marwadi for to blaze, then the Marwadi for sweat.
\end{quote}

\subsection*{You'll want to know: \mw{सीळो}}
\textbf{\mw{सीळो}} — \emph{sīḷo} — \textquotedblleft{}cool; pleasantly cold\textquotedblright{}.

\begin{quote}
\mw{आथण} \mw{हवा} \mw{सीळी} \mw{है।}
\end{quote}

\begin{quote}
\emph{āthaṇ havā sīḷī hai.}
\end{quote}

\begin{quote}
In the evening the breeze is cool.
\end{quote}

\begin{cousinweb}
\textbf{\mw{सीळो}} (\emph{sīḷo}) is inherited from Sanskrit \emph{śītala}, \textquotedblleft{}cold\textquotedblright{}; CDIAL lists Hindi \emph{sīlā}, \textquotedblleft{}cool, moist\textquotedblright{}, and a Gujarati cousin. CDIAL gives no Marwari form in this entry (CDIAL 12487).
\end{cousinweb}

\begin{grammarlens}[title={What you've built}]
One more word for this chapter.
\end{grammarlens}

\subsection*{Your turn}
\begin{itemize}
\raggedright
  \item \emph{Say it:} \emph{sīḷo}
  \item \emph{Say it:} \emph{sīḷo}, once more
  \item \emph{Say it:} \emph{āthaṇ havā sīḷī hai.} — In the evening the breeze is cool.
  \item \emph{Recall:} say the Marwadi for to blaze, then the Marwadi for sweat, then say \emph{sīḷo} again
\end{itemize}

\begin{tcolorbox}[breakable,colback=teal!4,colframe=teal!35!black,title={Before you move on}]
What does \mw{सीळो} mean? (\textquotedblleft{}Cool; pleasantly cold\textquotedblright{}.) Say it once more.
\end{tcolorbox}
\section[\texorpdfstring{\mw{सूकणो} (sūkṇo)}{सूकणो (sūkṇo)}]{\mw{सूकणो} (sūkṇo) — to dry up; to wither}

\label{lesson:MW-C228-sukno}



\begin{quote}
Before the new one: say the Marwadi for sweat, then the Marwadi for cool.
\end{quote}

\subsection*{You'll want to know: \mw{सूकणो}}
\textbf{\mw{सूकणो}} — \emph{sūkṇo} — \textquotedblleft{}to dry up; to wither\textquotedblright{}.

\begin{quote}
\mw{जेठ} \mw{मांय} \mw{खेत} \mw{सूकै} \mw{है।}
\end{quote}

\begin{quote}
\emph{jeṭh mā̃y khet sūkai hai.}
\end{quote}

\begin{quote}
The field dries up in early summer.
\end{quote}

Past (it …): \mw{वो} \mw{खेत} \mw{सूक्यो}, \emph{vo khet sūkyo}.

\begin{cousinweb}
\textbf{\mw{सूकणो}} (\emph{sūkṇo}) is inherited from a reconstructed \emph{*śuṣkati}, \textquotedblleft{}becomes dry\textquotedblright{}; CDIAL lists the Old Marwari verb \emph{sūkaï}, \textquotedblleft{}dries up\textquotedblright{} (CDIAL 12552). \textbf{\mw{सूको}}, dry, is its relative.
\end{cousinweb}

\begin{grammarlens}[title={What you've built}]
One more word for this chapter.
\end{grammarlens}

\subsection*{Your turn}
\begin{itemize}
\raggedright
  \item \emph{Say it:} \emph{sūkṇo}
  \item \emph{Say it:} \emph{sūkṇo}, once more
  \item \emph{Say it:} \emph{jeṭh mā̃y khet sūkai hai.} — The field dries up in early summer.
  \item \emph{Recall:} say the Marwadi for sweat, then the Marwadi for cool, then say \emph{sūkṇo} again
\end{itemize}

\begin{tcolorbox}[breakable,colback=teal!4,colframe=teal!35!black,title={Before you move on}]
What does \mw{सूकणो} mean? (\textquotedblleft{}To dry up; to wither\textquotedblright{}.) Say it once more.
\end{tcolorbox}
\section[(dialogue)]{Looking back: from Off to the Fair to Summer Sun}

\label{lesson:MW-R228-review-223-228}



\begin{quote}
Say the Marwadi for to swing and a big swing.
\end{quote}

\subsection*{Your turn}
Say these aloud:

\begin{itemize}
\raggedright
  \item to swing, a big swing, a big drum, a flag, a show
  \item to wander about, a puppet, a laddu, to leap up, a kettledrum
  \item a bundle, a cloth bag, a water-skin, bedding, a woollen blanket
  \item a sand hill, high, lonely, dry land, laid waste
  \item a well, deep down, thirsty, a Persian wheel, a drought
  \item parched, to blaze, sweat, cool, to dry up
\end{itemize}

\begin{tcolorbox}[breakable,colback=teal!4,colframe=teal!35!black,title={Before you move on}]
To swing? (\textbf{\mw{झूलणो}}, \emph{jhūlṇo}.) A big swing? (\textbf{\mw{हिंडोळो}}, \emph{hiṇḍoḷo}.) A big drum? (\textbf{\mw{ढोल}}, \emph{ḍhol}.) A flag? (\textbf{\mw{झंडो}}, \emph{jhaṇḍo}.) A show? (\textbf{\mw{तमासो}}, \emph{tamāso}.) To wander about? (\textbf{\mw{फिरणो}}, \emph{phirṇo}.) A puppet? (\textbf{\mw{कठपूतळी}}, \emph{kaṭhpūtḷī}.) A laddu? (\textbf{\mw{लाडू}}, \emph{lāḍū}.) To leap up? (\textbf{\mw{उछळणो}}, \emph{uchhaḷṇo}.) A kettledrum? (\textbf{\mw{नगाड़ो}}, \emph{nagāṛo}.) A bundle? (\textbf{\mw{गठड़ी}}, \emph{gaṭhṛī}.) A cloth bag? (\textbf{\mw{झोळी}}, \emph{jhoḷī}.) A water-skin? (\textbf{\mw{मसक}}, \emph{masak}.) Bedding? (\textbf{\mw{बिस्तर}}, \emph{bistar}.) A woollen blanket? (\textbf{\mw{कामळी}}, \emph{kāmḷī}.) A sand hill? (\textbf{\mw{टीबो}}, \emph{ṭībo}.) High? (\textbf{\mw{ऊंचो}}, \emph{ū̃cho}.) Lonely? (\textbf{\mw{सूनो}}, \emph{sūno}.) Dry land? (\textbf{\mw{थळ}}, \emph{thaḷ}.) Laid waste? (\textbf{\mw{उजाड़}}, \emph{ujāṛ}.) A well? (\textbf{\mw{कूवो}}, \emph{kūvo}.) Deep down? (\textbf{\mw{ऊंडो}}, \emph{ū̃ḍo}.) Thirsty? (\textbf{\mw{तिसायो}}, \emph{tisāyo}.) A Persian wheel? (\textbf{\mw{रहट}}, \emph{rahaṭ}.) A drought? (\textbf{\mw{दुकाळ}}, \emph{dukāḷ}.) Parched? (\textbf{\mw{लूखो}}, \emph{lūkho}.) To blaze? (\textbf{\mw{तपणो}}, \emph{tapṇo}.) Sweat? (\textbf{\mw{पसीनो}}, \emph{pasīno}.) Cool? (\textbf{\mw{सीळो}}, \emph{sīḷo}.) To dry up? (\textbf{\mw{सूकणो}}, \emph{sūkṇo}.)
\end{tcolorbox}
